\section{Objetivos}

\subsection{Objetivo general}

Contribuir a la comprensión de los factores que hacen que las estimaciones de tiempo de viaje de Google Maps se alejen del tiempo real en los desplazamientos de los estudiantes de la UPTC por la Avenida Oriental de Tunja, como base para que estos factores puedan anticiparse y estudiarse en investigaciones posteriores.

\subsection{Objetivos específicos}

\begin{enumerate}
    \item Determinar la diferencia entre el tiempo estimado por Google Maps y el tiempo real de recorrido en el corredor de la Avenida Oriental, en ambos sentidos.

    \textbf{Meta:} reportar la diferencia promedio (en minutos y en porcentaje) y el error absoluto medio, por sentido, con los trayectos medidos durante el periodo de recolección.

    \item Comparar la diferencia entre el tiempo estimado y el tiempo real en franjas de hora pico y no pico.

    \textbf{Meta:} reportar la diferencia promedio y su variabilidad por franja horaria y por sentido.

    \item Identificar si los eventos o cierres viales registrados durante las mediciones se asocian con diferencias mayores entre el tiempo estimado y el real.

    \textbf{Meta:} reportar la diferencia promedio de los trayectos con y sin evento o cierre vial registrado, indicando cuántos trayectos hubo en cada grupo.

    \item Identificar cuáles de las variables estudiadas (franja horaria, sentido del recorrido y eventos o cierres viales) se asocian en mayor medida con la diferencia entre el tiempo estimado y el real.

    \textbf{Meta:} reportar, para cada variable, la magnitud de su asociación con la diferencia de tiempos y ordenarlas según su importancia.
\end{enumerate}